%%
%% appendix_benchmark_methodology.tex
%% Online Supplementary Material for:
%%   "Edge-Intelligent Video Analytics: A Systematic Survey of
%%    Architectures, Applications, and Deployment Strategies
%%    Across Smart City Verticals"
%% ACM Computing Surveys — Supplementary Appendix
%%
%% This document is published as online-only supplementary material
%% in the ACM Digital Library alongside the main article.
%% Compile standalone: pdflatex appendix_benchmark_methodology.tex
%%

\documentclass[manuscript]{acmart}
\usepackage{booktabs}
\usepackage{enumitem}
\setcopyright{none}

\title{Supplementary Material: EdgeVA Benchmarking Methodology}
\subtitle{Online Appendix for the ACM Computing Surveys Submission}

\author{Mukhiddin Toshpulatov}
\orcid{0000-0002-8819-852X}
\affiliation{%
  \institution{Computer Engineering, College of Information Technology
               Convergence, Gachon University}
  \city{Seongnam}
  \country{South Korea}}
\email{mukhiddin@gachon.ac.kr}

\author{Wookey Lee}
\orcid{0000-0001-8586-4577}
\authornote{Corresponding authors.}
\affiliation{%
  \institution{Departments of Industrial and Biomedical Science Engineering,
               Inha University}
  \city{Incheon}
  \country{South Korea}}
\email{trinity@inha.ac.kr}

\author{Jinsoo Cho}
\orcid{0000-0002-8119-8068}
\authornotemark[1]
\affiliation{%
  \institution{Computer Engineering, College of Information Technology
               Convergence, Gachon University}
  \city{Seongnam}
  \country{South Korea}}
\email{jscho@gachon.ac.kr}

\renewcommand{\shortauthors}{Toshpulatov, Lee, and Cho}

\begin{abstract}
This supplementary appendix provides the complete EdgeVA benchmarking
methodology supporting the hardware performance results in the main
article. It documents hardware configurations, software versions, model
export procedures, measurement protocols, pipeline-timing breakdowns,
and confidence-interval calculations. All scripts and raw results are
publicly available at \url{https://github.com/muxiddin19/EdgeVA}.
\end{abstract}

\begin{document}
\maketitle

\section{Hardware Configuration}

\subsection*{Server Platform (CPU and GPU Runs)}

\begin{itemize}[leftmargin=*, noitemsep]
  \item \textbf{CPU}: Intel Core i9-14900K (24-core/32-thread,
    3.2\,GHz base / 6.0\,GHz boost)
  \item \textbf{GPU}: Dual NVIDIA GeForce RTX~4090 (24\,GB GDDR6X each);
    benchmarks pinned to GPU~0 (\texttt{device\_id=0})
  \item \textbf{RAM}: 125\,GB DDR5
  \item \textbf{OS}: Ubuntu~22.04.4 LTS (kernel 6.8.0)
  \item \textbf{Driver / CUDA}: NVIDIA driver 535.288.01; CUDA~12.2
  \item \textbf{Environment}: dedicated \texttt{edgeva} conda environment
    (Python~3.10.20), isolated from other projects to avoid library conflicts
  \item \textbf{Key packages}: \texttt{onnxruntime-gpu}~1.24.4 (PyPI),
    \texttt{ultralytics}~8.4.36, \texttt{onnxslim}~0.1.90,
    \texttt{onnx}~1.16.0, \texttt{scipy}~1.13.0, \texttt{numpy}~1.26.4
\end{itemize}

\subsection*{Jetson Orin NX Figures (from Literature)}

Jetson Orin NX 16\,GB results cited from~\cite{alqahtani2024edgebench}:
JetPack~5.1.2, TensorRT~8.6.2, ONNX~1.14, INT8 calibration on
COCO 2017 val (500 representative images selected by entropy).

\section{Model Export Procedure}

ONNX models exported using the \texttt{ultralytics} Python API:

\begin{verbatim}
from ultralytics import YOLO
model = YOLO("yolov8n.pt")
model.export(format="onnx", opset=20,
             simplify=True, dynamic=False,
             imgsz=640)
\end{verbatim}

Post-processed with \texttt{onnxslim} for graph simplification.
No TensorRT engine conversion was performed for server-side runs;
the ONNX Runtime \texttt{CUDAExecutionProvider} with
\texttt{ORT\_ENABLE\_ALL} graph optimization was used exclusively.

\subsection*{Model Checksums (SHA-256, truncated)}

\begin{tabular}{lll}
\toprule
\textbf{Model} & \textbf{SHA-256 (first 8 chars)} & \textbf{Size} \\
\midrule
\texttt{yolov8n.onnx} & \texttt{b3a1f2e4} & 12.3\,MB \\
\texttt{yolov8s.onnx} & \texttt{d7f2c8a1} & 42.8\,MB \\
\texttt{yolov8m.onnx} & \texttt{e9c4b3f7} & 99.0\,MB \\
\bottomrule
\end{tabular}

Full SHA-256 hashes and download links are in
\texttt{experiments/results/gpu\_summary.json} in the repository.

\section{Measurement Protocol}

\begin{itemize}[leftmargin=*, noitemsep]
  \item \textbf{Input}: random FP32 tensor $1{\times}3{\times}640{\times}640$,
    values $\sim\mathcal{U}(0,1)$; batch size~=~1.
  \item \textbf{Warmup}: 100 inference calls discarded before timing begins.
  \item \textbf{Timed runs}: 500 calls; wall-clock time measured with
    \texttt{time.perf\_counter()} bracketing each \texttt{session.run()} call.
  \item \textbf{Reported statistics}: p50 (median), p95, p99 latency (ms);
    FPS $= 1000 / \mathrm{p50}$; mean and std also recorded.
  \item \textbf{Threading}: \texttt{intra\_op\_num\_threads=8} for CPU
    provider; GPU provider uses the CUDA default stream.
  \item \textbf{Thermal steady-state}: server fan control set to automatic;
    10-minute system idle before each benchmark run.
    GPU junction temperature remained below 83°C throughout.
  \item \textbf{Power measurement}: GPU power sampled via
    \texttt{nvidia-smi --query-gpu=power.draw --format=csv}
    at 1\,s intervals during the 500-run window; mean reported.
    GPU~0 drew $\approx$68\,W (YOLOv8n CUDA) to
    $\approx$112\,W (YOLOv8m CUDA) during inference.
  \item \textbf{CPU provider isolation}: CPU runs used
    \texttt{CPUExecutionProvider} only; GPU was idle.
\end{itemize}

\section{Pipeline-Timing Breakdown}

Full-pipeline timing (reported in the main article, Table~2 and
Section~6.4) applies the same 500-run / 100-warmup protocol to a
sequential five-stage pipeline:

\begin{enumerate}[leftmargin=*, noitemsep]
  \item \textbf{Preprocess}: BGR$\rightarrow$RGB, HWC$\rightarrow$BCHW,
    divide by 255, cast to FP32 (\texttt{numpy}).
  \item \textbf{Inference}: ONNX Runtime \texttt{session.run()}.
  \item \textbf{Postprocess}: confidence threshold filter ($>0.25$),
    box decoding.
  \item \textbf{Tracking}: LITE-proxy association (Hungarian algorithm,
    IoU cost matrix, $\tau_\mathrm{IoU}=0.45$).
  \item \textbf{Analytics}: zone occupancy count (point-in-polygon test).
\end{enumerate}

Each stage is timed separately with \texttt{time.perf\_counter()};
total pipeline latency = sum of all five stages.

\section{Confidence Intervals}

With 500 runs and empirical latency standard deviations of
2--5\,ms (CUDA) and 3--12\,ms (CPU), the 95\% bootstrap confidence
interval on the p50 estimate is $\pm0.3$--$0.7$\,ms
($B=10{,}000$ bootstrap replicates). Full bootstrap distributions
are provided in \texttt{experiments/results/gpu\_summary.json}.

The CI formula used:

\begin{equation}
\mathrm{CI}_{95} = \left[
  \hat{\theta}^*_{0.025},\;\hat{\theta}^*_{0.975}
\right],
\end{equation}

where $\hat{\theta}^*_\alpha$ is the $\alpha$-quantile of the
bootstrap distribution of the median estimator over $B$ resamples.

\section{Reproducibility Checklist}

\begin{tabular}{lc}
\toprule
\textbf{Item} & \textbf{Available} \\
\midrule
Benchmark scripts (\texttt{run\_benchmarks\_gpu.py}) & \checkmark \\
Model export script (\texttt{export\_models.py}) & \checkmark \\
Environment spec (\texttt{setup\_env.sh}, \texttt{requirements.txt}) & \checkmark \\
Raw CSV results (\texttt{experiments/results/}) & \checkmark \\
Bootstrap summary JSON (\texttt{gpu\_summary.json}) & \checkmark \\
Model ONNX files (via HuggingFace or direct download) & \checkmark \\
\bottomrule
\end{tabular}

\medskip
\noindent
All artefacts are at:
\url{https://github.com/muxiddin19/EdgeVA}

\bibliographystyle{ACM-Reference-Format}
\bibliography{../../AntVision/csur_submission/AntVision_refs,
              ../../AntVision/csur_submission/csur_extra}

\end{document}
